\begin{figure}[ht]
\centering
\begin{tikzpicture}[x=0.34cm, y=0.34cm, font=\footnotesize]

% Correct decimal digits vs iteration for fixed-point (linear),
% Newton (quadratic), and Halley (cubic) on the same simple root.
\draw[->, thick] (0,0) -- (13,0) node[right] {iteration $k$};
\draw[->, thick] (0,0) -- (0,18) node[above, align=center]
  {$-\log_{10}\abs{e_{k}}$\\(correct digits)};

\foreach \k in {0,2,4,6,8,10,12} {
  \draw (\k,0) -- (\k,-0.3);
  \node[anchor=north, font=\scriptsize] at (\k,-0.3) {\k};
}
\foreach \d in {0,4,8,12,16} {
  \draw (0,\d) -- (-0.3,\d);
  \node[anchor=east, font=\scriptsize] at (-0.3,\d) {\d};
}

% Fixed-point: linear, ~0.8 digits per step
\draw[thick, engblue] plot coordinates {
  (0,0.5) (2,2.1) (4,3.7) (6,5.3) (8,6.9) (10,8.5) (12,10.1)
};
\node[engblue, anchor=west, font=\scriptsize] at (12.1,10.1) {linear};

% Newton: quadratic, digits roughly double
\draw[thick, engred] plot[smooth] coordinates {
  (0,0.5) (1,1.1) (2,2.2) (3,4.4) (4,8.8) (5,16) (8,16) (12,16)
};
\node[engred, anchor=west, font=\scriptsize] at (5.2,15.2) {Newton (order 2)};

% Halley: cubic, digits roughly triple
\draw[thick, engamber] plot[smooth] coordinates {
  (0,0.5) (1,1.6) (2,4.8) (3,14.4) (4,16) (8,16) (12,16)
};
\node[engamber!70!black, anchor=west, font=\scriptsize] at (3.6,13.0)
  {Halley (order 3)};

% Double-precision ceiling
\draw[thick, dashed, enggray] (0,16) -- (13,16);
\node[enggray, anchor=west, font=\scriptsize] at (5.0,16.8)
  {binary64 ceiling};

\end{tikzpicture}
\caption[Convergence of fixed-point, Newton, and Halley iterations]{Correct
decimal digits against iteration count for a linearly convergent fixed-point
map, Newton's quadratic iteration, and Halley's cubic iteration on the same
simple root. The linear map adds a fixed number of digits per step; Newton
roughly doubles the digit count; Halley roughly triples it, reaching the
binary64 ceiling in three or four steps at the cost of one extra derivative
evaluation per step. The higher order buys fewer iterations only when each
evaluation is cheap relative to the setup, which is the trade the engineer
weighs before reaching past Newton. Source: authors' computation; the slopes
are the theoretical convergence orders.}
\label{fig:vol02ch16:halley-newton}
\end{figure}
